\section{Síntese}

\begin{frame}{Perguntas para avaliar uma infraestrutura}
  \begin{enumerate}
    \item Sabemos de quais componentes e fornecedores dependemos?
    \item Conseguimos auditar e operar os componentes críticos?
    \item Temos competências para corrigir e manter os sistemas?
    \item Existem alternativas e caminhos de migração testados?
    \item Quem define as regras e responde por decisões e incidentes?
  \end{enumerate}
  \note{Retome o exemplo da abertura e aplique estas perguntas. Deixe claro que se trata de uma heurística de análise, não de uma certificação formal.}
\end{frame}

\begin{frame}{Conclusões}
  \begin{itemize}
    \item Soberania digital é multidimensional.
    \item Dependências atravessam camadas técnicas e institucionais.
    \item Autonomia requer capacidades concretas, alternativas e governança.
  \end{itemize}
  \vspace{1em}
  \begin{block}{Pergunta para discussão}
    Quais componentes críticos da infraestrutura que utilizamos conseguimos efetivamente operar e substituir?
  \end{block}
  \note{Ajuste estas conclusões ao conjunto de evidências que de fato aparecerá na versão final. Não apresente como resultados demonstrados aquilo que for apenas hipótese.}
\end{frame}
